\sectionkicker{Theme synthesis}
\subsection*{横断比較 — Tool-Native化はmodel、API、runtimeのどこで起きたか}
\addcontentsline{toc}{subsection}{横断比較 — Tool-Native化はmodel、API、runtimeのどこで起きたか}
Qwen、GLM、SGLangを同じmodel比較へ押し込むのではなく、tool-enabled API、model release、runtime uptakeという層に分けると、1月にstackの上下が同時に変わったことが見える。
\medskip
\begin{center}
\small
\renewcommand{\arraystretch}{1.16}
\begin{tabularx}{\linewidth}{@{}p{0.20\linewidth}X@{}}
\toprule
観察軸 & 1月の一次資料・論文から読めること \\
\midrule
reasoning mode & Alibaba Cloud Model Studioはqwen3-max-2026-01-23についてthinking/non-thinking modeの統合を説明する。model lifecycleだけでなく、同一managed endpointでreasoning modeをどう切り替えるかがAPI surfaceになる。 \cite{src-eac8dfbeb78c} \\
tools during thinking & 同更新はthinking中のweb search、web extraction、code interpreter利用を記録する。tool useが外付けharnessだけでなくmanaged model/APIのreasoning pathへ組み込まれている。 \cite{src-eac8dfbeb78c} \\
broader availability & 同じlifecycle pageには1月4日、tool-oriented／codingやmultimodal modelを含む複数のQwen3・Wan 2.6系のglobal availabilityも記録される。単一model releaseより広いmanaged catalog更新が同月に進んだ。 \cite{src-eac8dfbeb78c} \\
region / living docs & Model Studioのavailability dateはregion差を持ち得て、lifecycle page自体もliving documentationである。January chronologyはcurrent default aliasではなくdated entryとdeployment regionへ固定する。 \cite{src-eac8dfbeb78c} \\
GLM-Image & Z.aiは1月14日のGLM-Imageをautoregressive semantic understandingとdiffusion-based decodingを組み合わせるcontrollable image-generation modelとして説明した。multimodal生成側のreleaseである。 \cite{src-9156c37fb1c2} \\
GLM-4.7-Flash & 1月19日のGLM-4.7-Flashはcoding・reasoning・generative task向けのlightweight／free-tier variantとして記録される。GLM-Imageとは別release・別役割であり、一つのarchitectureとして混同しない。 \cite{src-9156c37fb1c2} \\
vendor claim boundary & Qwen lifecycle catalogやZ.ai release notesのcapability、quality、latency、throughput、best-in-class等の比較表現はvendor-reportedである。release availabilityの一次事実と性能評価を分ける。 \cite{src-eac8dfbeb78c,src-9156c37fb1c2} \\
runtime uptake & SGLang v0.5.8はGLM-4.7-Flashのday-zero supportを含むmodel/runtime更新をまとめた。これはimplementation availabilityのEvidenceであり、同等品質やproduction validationを意味せず、release内speedup値もproject claimである。 \cite{src-9156c37fb1c2,src-6a97353d35a9} \\
\bottomrule
\end{tabularx}
\renewcommand{\arraystretch}{1.0}
\end{center}
\normalsize
\begin{claimboundary}[この比較の境界]
この表は本号で既に検証・参照している一次資料と論文を横断比較の形へ再配置したものである。vendor / project / author が報告した性能値や優位性は独立再現済みの一般則へ変換せず、本文とSource Notesの帰属境界をそのまま維持する。
\end{claimboundary}
